\newpage
\section{The Proposed System}

\subsection{Business Context}

\subsubsection{Application domain}

The application domain is career guidance and exploration, situated at the intersection of educational technology and occupational information provision. Participants in this domain are individuals facing an occupational decision, institutions supporting them, and the bodies that publish occupational information. The domain's characteristic activity is the transfer of occupational information to a person who must act on it, and Chapter 1 established that this transfer is currently incomplete: guidance provision is dominated by instruments that elicit self-\allowbreak{}description rather than convey what work involves.

The proposed system, JobQuest, addresses this by allowing a user to perform representative tasks of an occupation, receive an assessment of the competencies those tasks elicited, and use that assessment to determine which further occupations to explore. It is a platform for informed exploration rather than a placement, recruitment, or credentialing service, and this boundary determines several of the business rules below.

\subsubsection{Business processes}

Three processes constitute the operation of the platform. They are modelled formally in Section 5.1.1.

The exploration process is the primary user-facing process. A registered user optionally completes an orientation conversation, which is provided for users who do not know which occupation to begin with and is not a prerequisite. The user selects an occupation and a level within it, is shown an introduction to the role and team, may open the full working context before any task begins, and then works through the tasks comprising that level. Each task is a scenario drawn from the occupation and is answered in one of nine formats --- multiple choice, multiple selection, ordering, matching, true/false, timed choice, free text, simplified email, or simplified code --- with hints available throughout. Tasks ask what the user does in the situation, a behavioural-\allowbreak{}tendency instruction\cite{r28}, because how a user tends to act is itself information about fit; scoring uses expert effectiveness keys in either case. Each response option is of one of three types: correct, the most effective response; trade-off, effective for one competency at a cost to another; or wrong, ineffective or harmful. Options and rubric criteria are scored on a single ordinal scale per competency, from $-$1 (ineffective or harmful) through 0 and +1 to +2 (most effective), so that evidence from choices and from written responses accumulates in the same competence record and agreement with experts can be computed on an ordinal scale. Using a hint is recorded and lowers the weight of that response as evidence of independent competence without blocking progress; a timed task that expires is recorded as unanswered, with its response time, and carries no separate penalty. The system evaluates each response, records the competencies it evidences, and returns a result, feedback specific to the response, and guidance for improvement. Completing a level awards competence credit; accumulated credit unlocks higher levels of the same occupation or, where competencies overlap sufficiently, entry to an adjacent occupation above its entry level. On completing an occupation, the user may rate it on five attributes, which contribute to that occupation's aggregate profile.

The content lifecycle process governs how occupations enter the catalogue, and it involves two roles with distinct responsibilities. Domain experts own the occupational content and author it in the system: for each scenario they enter the required facts, the competencies it must elicit, the critical incidents from which response options are drawn, the effectiveness of each option, and the rubric against which free responses are judged, which together form the scenario's seed. A seed is a template rather than a script: it fixes the canon events of a scenario --- the situations every run must pass through, at set points --- together with the keyed options and rubric attached to them, and leaves the narrative between those events to be developed by the model in response to the user's answers. Validation accordingly has two layers. In the first, the system generates a sample of scenarios from the seed, and experts other than the seed's authors review the seed and the sample in the system, judging their accuracy and returning approvals, rejections, or corrections, each recorded with the identity of the reviewing expert. In the second, every scenario generated at run time is checked automatically against its seed --- that it passes through every canon event, preserves the required facts, and leaves each option's keyed profile unchanged (NFR-15) --- and one that fails is regenerated or replaced by an instance from the approved sample. The content administrator does not author occupational content. The administrator manages the occupation as an entry in the catalogue, arranges how approved content is laid out and presented in gameplay, reviews the occupation as a whole before release, publishes it, and operates its lifecycle thereafter. Published content is revalidated annually by the experts, since occupational practice changes. In the present phase, in which no industry experts are yet engaged, the project team performs the expert role on provisional seeds (Section 4.2.1).

The evaluation process is internal and supports the research objectives. It proceeds in two stages (Section 4.2.2): a synthetic benchmark constructed by the project team, and a reference set of real responses and generated scenarios rated independently by domain experts, against which generation quality and assessment accuracy are measured offline as specified in Section 2.6.

\subsubsection{Business rules}

Business rules are stated here as policies and constraints of the domain that hold independently of any particular implementation. Constraints arising from implementation choices --- which models are used, how cost is bounded, where data is processed --- are recorded separately as design constraints in Section 4.3.3.

{\small\setstretch{1}\renewcommand{\arraystretch}{1.1}
\begin{longtable}{|>{\raggedright\arraybackslash}p{\dimexpr 0.103\textwidth-2\tabcolsep-\arrayrulewidth\relax}|>{\raggedright\arraybackslash}p{\dimexpr 0.597\textwidth-2\tabcolsep-\arrayrulewidth\relax}|>{\raggedright\arraybackslash}p{\dimexpr 0.295\textwidth-2\tabcolsep-\arrayrulewidth\relax}|}
\caption{Business rules governing the proposed system.} \\
\hline
\textbf{ID} & \textbf{Rule} & \textbf{Origin} \\ \hline
\endfirsthead
\hline
\textbf{ID} & \textbf{Rule} & \textbf{Origin} \\ \hline
\endhead
BR-01 & The entry level of every occupation is accessible without prerequisite. & Exploration is the purpose; gating discovery defeats it \\ \hline
BR-02 & Access to a higher level requires the competence credit specified for that level. & Competency-\allowbreak{}based progression (Section 2.2) \\ \hline
BR-03 & Demonstrated competencies that overlap an adjacent occupation's requirements, as determined from the occupational taxonomy, count towards that occupation's levels, so that a user may enter it above its entry level. & Sections 2.2, 2.6; a user exploring related work need not restart \\ \hline
BR-04 & An occupation may be rated only by a user who has completed it. & Rating validity depends on exposure \\ \hline
BR-05 & No scenario seed or rubric is published without domain-\allowbreak{}expert validation of the seed and of a sample of the scenarios generated from it. & Section 2.5; accuracy of occupational representation \\ \hline
BR-06 & Assessment results are formative. They are not transmitted to employers and do not determine access to any external opportunity. & Section 2.7; proportionality to assessment reliability \\ \hline
BR-07 & Occupational representation must include the unfavourable aspects of the work. & Realistic job preview evidence (Section 2.1) \\ \hline
BR-11 & Replaying a scenario cannot raise a competency beyond the credit that scenario can award. & Keeps unlocks tied to demonstrated competence (BR-02) \\ \hline
BR-13 & An occupation's content becomes available to explorers only when the content administrator publishes it, which requires approved seeds (BR-05). & Scenario lifecycle (Section 5.1.3) \\ \hline
BR-14 & Revising a seed withdraws the scenarios generated from it until the revised seed is approved. & Scenario lifecycle; accuracy of occupational representation \\ \hline
BR-15 & Published content is revalidated by domain experts every year. & Occupational practice changes \\ \hline
BR-16 & A retired occupation remains in the records of users who attempted it but can no longer be started. & Scenario lifecycle; continuity of users' records \\ \hline
\end{longtable}
}

BR-01 and BR-03 are complementary. BR-01 keeps discovery free: the entry level of any occupation can be tried without prerequisite. BR-03 concerns what earlier work is worth: competence already demonstrated in one occupation is credited towards an adjacent one, so that the user does not begin it again from the entry level.

\subsubsection{Organisational policies and legal-\allowbreak{}ethical considerations}

The platform processes personal data, including performance data, and is therefore subject to Vietnam's Law on Personal Data Protection and its implementing decree\cite{r99,r100} and to comparable frameworks for users in other jurisdictions. Three policies follow. Data minimisation limits collection to what the exploration and assessment functions require. Aggregate-\allowbreak{}only publication governs any statistic derived from user performance, with suppression thresholds applied so that small populations cannot be de-\allowbreak{}anonymised; as a privacy measure rather than a rule of the domain, it is specified as a non-\allowbreak{}functional requirement (NFR-06). Formative-\allowbreak{}use limitation (BR-06) prevents assessment output from acquiring consequences the method's reliability does not support.

One provision of the law shapes the design directly. Sending a user's response to a model provider located abroad is a cross-\allowbreak{}border transfer of personal data; responses are therefore stripped of identifying details before they leave the platform, since de-\allowbreak{}identified data is not personal data under the law. How this provision applies to a university-\allowbreak{}hosted project, and whether behavioural logs fall within the decree's sensitive categories, are to be confirmed with the university's legal office.

Two further policies concern the system's use of generative models. Provenance disclosure requires that users be informed that scenarios are machine-\allowbreak{}generated from expert-\allowbreak{}authored material and that assessments are machine-\allowbreak{}produced. Expert accountability assigns responsibility for the occupational accuracy of published content to the authoring and reviewing domain experts, as recorded in the system, rather than to the generating model. Under the intended partnership described in Section 4.1.5, data governance would additionally follow the requirements of the Ministry of Education and Training for systems serving students. Detailed processing policies, including third-party model providers, are specified in Section 5.3.10.

\subsubsection{Business model and organisation}

JobQuest is designed as a public-\allowbreak{}benefit platform rather than a commercial product, and its business model follows from that position and from the evidence on which the system rests.

Value and beneficiaries. The platform offers learners experiential occupational information before commitment, and offers education authorities and institutions a scalable complement to limited counselling capacity. All content is available to all users without charge; access to levels and occupations is governed by demonstrated competence (BR-01 to BR-03) rather than by payment. The choice is deliberate: the users who most need occupational information --- secondary students choosing a field under examination and family pressure --- are those least able to pay for it, and the platform's effect depends on reaching them.

Partners and funders. The Ministry of Education and Training and educational-\allowbreak{}technology organisations may each act both as partners and as funders. As a partner, the Ministry --- preceded by city-level education departments --- supplies distribution through schools, legitimacy with schools and families, and a governance framework for students' data, while edtech organisations supply pilot hosting, distribution, and technical cooperation. As funders, both may provide programme funding, grants, or in-kind support such as cloud credits. Industry partners and professional associations supply and coordinate domain experts, and the host university provides academic oversight. One set of safeguards applies to every partner and funder, whatever its role: none authors, validates, or ranks occupational content; none receives identifiable user data, only aggregate statistics (NFR-06); none may introduce charges, promotional content, or monetisation of users; and every partnership and source of funding is disclosed on the platform.

Funding and sustainability. Without user revenue, operation is funded through public and philanthropic support: grant funding sought from the Ministry and from educational-\allowbreak{}technology organisations; sponsorship from companies, which may contribute experts as well as funds; voluntary donations through a donation box on the platform; and in-kind contribution of expert consultation. The first expert panel, in Stage 1, contributes voluntarily; from Stage 2 an honorarium for experts is budgeted within grant funding. The cost structure determines how that funding must be managed. The dominant variable cost is language-\allowbreak{}model inference, incurred per session, which is why per-session cost is a monitored requirement (NFR-11). The system manages it by enforcing a configured ceiling on the model cost of each session, a cost-\allowbreak{}management mechanism specified with the design in Chapter 5. The dominant fixed cost is expert authoring, incurred per occupation, which is why coverage expands field by field rather than all at once. Model access is provider-\allowbreak{}portable (NFR-12, DC-02), so that inference can move to lower-cost or institutionally hosted models as funding conditions change. Two further levers reduce cost. A panel of smaller judging models from different families is less expensive than a single large judge\cite{r74}, and because a seed fixes a scenario's canon events, keyed options, and rubric in advance (BR-05), run-time generation is confined to the narrative between those events, which bounds the length and therefore the cost of each generation.

Sponsorship and content integrity. A model in which employers paid to shape how occupations are portrayed would reward favourable portrayal, in direct conflict with the realistic job preview evidence (Section 2.1) and with BR-07. Sponsorship is therefore admitted only as acknowledgement. Sponsoring companies are named, with their logos, under a plain ``Supported by'' heading; the acknowledgement carries no slogans, product claims, or links to product pages, and it never appears within scenarios, assessments, or rankings. Sponsors do not author, approve, or rank content, and every sponsorship is disclosed. Whether such an acknowledgement falls within the advertising regulations is subject to legal review, to be settled with the university's legal or partnership office before any logo is displayed.

Perception of a free service. A service offered without charge can be taken to be of low quality. The platform answers this with signals of credibility rather than with a price. Sponsor acknowledgements and the donation box are present from the outset. A badge marks scenarios validated by domain experts; it is shown only on scenarios that experts have actually approved, and therefore from Stage 1. Evaluation results are published once Stage A or Stage B has produced them, with the stage stated, since Stage A does not establish agreement with experts. Endorsement by the Ministry or by partner institutions is deferred to Stage 3. No paid tier is introduced.

Starting point and growth roadmap. A public authority does not adopt an unproven platform, so the project begins where it stands and advances only as evidence accumulates. The policy direction supports the end goal: the national scheme on career guidance and student streaming in general education, approved by Decision 522/QĐ-TTg, covered the period 2018--2025\cite{r95}; its results were reviewed in November 2025, with a decree on career guidance and educational streaming in preparation\cite{r96}; and Ho Chi Minh City has set a 2030 target of at least 35\% of students in basic sciences, engineering, and technology\cite{r97}. Table 4.2 sets out the stages. Each has a gate --- the evidence required before the next begins --- and meeting the thresholds of the two research problems is the evidence that justifies approaching funders and authorities.

The stages, their periods, activities, funding, and the gate to the next stage are tabulated in Appendix 3.

The project remains university-\allowbreak{}hosted during Stages 0 and 1, which gives it institutional standing without a separate legal entity. The operating organisation formed in Stage 2 is to take a non-\allowbreak{}commercial form consistent with free access, such as a social enterprise or a university-\allowbreak{}affiliated centre, to be settled with legal advice at that stage.

Organisational structure. The organisation is staged. During the present phase and the capstone phase, the three-\allowbreak{}member project team holds every operating role, one member holding several where necessary; domain experts join from the capstone phase; and the project advisor provides academic oversight. For sustained operation after the capstone phase, the proposed structure places a steering committee --- with representatives of the Ministry, the host university, funding organisations, and the project lead --- above five operating units, advised by an expert board of senior domain and assessment specialists. Table 4.3 sets out the units, their positions, and their responsibilities. The indicative core is about eleven full-time positions, with domain leads part-time and contributing experts engaged per occupation, a size chosen to remain sustainable on grant funding. The numbers of positions are indicative and the structure is not rigid: in the early stages one person may hold more than one position.

The units, positions, and responsibilities are tabulated in Appendix 3.

Occupational domains are the unit of content governance: each has a domain lead accountable for its accuracy. Within the information technology field covered in the present phase, the proposed domains are software engineering; data and artificial intelligence; network and systems infrastructure; information security; and IT product and business analysis. Each further field is added as a new domain with its own lead, which is how coverage expands without restructuring the organisation.

\subsection{System Description}

\subsubsection{System overview}

JobQuest is a web-based platform that presents occupations as sequences of performable tasks, assesses the competencies a user demonstrates while performing them, and uses those assessments to structure further exploration. Its purpose is to supply occupational information experientially, at a breadth that authored content cannot reach, to people who must choose or change an occupation.

The experience is delivered through three presentation and progression conventions selected against the evidence in Section 2.3 rather than by preference. Progression follows a role-\allowbreak{}playing convention, in which levels and access to further occupations are earned through demonstrated competence rather than time spent, which the role-play evidence supports as the genre most associated with skill acquisition. Scenarios are presented in a branching narrative form comparable to a visual novel, the narrative moderator being among the strongest identified in the gamification meta-\allowbreak{}analyses. Career and life events are introduced probabilistically across sessions, a convention drawn from roguelike design, so that the circumstances bearing on occupational fit vary between runs; consistent with the harm evidence in Section 2.3, this randomisation governs narrative variety only and is not coupled to any reward schedule or monetisation. These conventions carry engagement; consistent with Section 2.8 they are not relied upon for competence development, which rests on task authenticity and assessment.

The system serves four user groups: the explorer, the domain expert, the content administrator, and the system administrator.

The explorer is the primary user: a student choosing a field of study, a graduate entering the labour market, or a worker considering transition. The explorer performs tasks, receives assessment and guidance, accumulates competence credit, and navigates among occupations. All other groups exist to make this possible.

The content administrator manages the occupation catalogue and is responsible for how expert-\allowbreak{}authored content reaches explorers: arranging its layout and presentation in gameplay; creating, reviewing, and publishing occupations; and operating the content lifecycle, including monitoring coverage and currency across occupations and arranging revalidation. The content administrator does not author or alter occupational content, which belongs to the domain experts.

The system administrator operates the platform: configuring and monitoring the generation and assessment models, managing capacity and cost, administering accounts and access, and overseeing data protection controls.

Domain experts are senior or lead practitioners in an occupation, typically with around five or more years in the role and experience supervising junior staff, who therefore know what entry-level work demands and where newcomers typically fail. They use the system directly: they author and own the occupational content of scenarios --- the required facts, target competencies, response options and their effectiveness keys, and rubrics --- and they review the scenarios the system generates and the cases in which automated judgment departs from their own. Each role is served by three to four experts, three being the accepted minimum for content validation\cite{r69}, drawn from at least two types of organisation --- such as product, outsourcing, startup, and in-house enterprise IT --- because the same job title involves different work in different organisations; no expert reviews scenarios generated from a seed they helped prepare. Their involvement is heaviest when an occupation's seeds are first prepared and lighter thereafter, through annual revalidation, and is arranged part-time through partner organisations; the first panel is recruited through the advisor's and the university's alumni networks and contributes voluntarily. Across the approximately twenty-two IT roles in scope, the full network is some sixty-six to eighty-\allowbreak{}eight experts, reached progressively (Section 4.1.5). In the present phase the project team prepares provisional seeds from authoritative public sources under the advisor's guidance, and these are not presented as expert-\allowbreak{}validated. This group is the guarantor of the system's central claim, that its representations correspond to real work.

\subsubsection{Research problems}

In this report, a research problem is a question that meets three conditions: the answer is not established in the existing literature; the answer constitutes knowledge that generalises beyond this particular system; and the project's evaluation produces evidence that answers it. A question resolved by a design decision fails the second condition, and a question about how evidence is gathered concerns method rather than a problem. Two questions meet all three conditions, both in artificial intelligence applied to education, and they correspond to the first two specific objectives in Section 1.2.

RP-1, grounded generation of occupational scenarios. The literature establishes that language models generate fluent instructional content but that fidelity to a specified domain requires explicit constraint (Section 2.5); it does not establish how accurate expert-\allowbreak{}seeded generation is for task-based occupational content, nor whether generated tasks preserve the competency they target (Section 2.2). The question is therefore: when generation is constrained by an expert-\allowbreak{}authored seed specifying required facts and target competencies, what proportion of generated scenarios do domain experts judge to be occupationally accurate and to require the target competency, and how does this compare with generation from the same instructions without the seed? The comparison is what makes the answer knowledge rather than description, since without it an accuracy figure cannot be attributed to the seeding.

RP-2, validity of language-\allowbreak{}model stealth assessment. Reported agreement between automated judges and human raters is systematically optimistic, and chance-\allowbreak{}corrected, difficulty-\allowbreak{}stratified validation is rarely reported (Section 2.5); nor has such assessment been validated where the evidence of competence is the performance of a task rather than an answer to a test item (Section 2.2). The question is therefore: to what extent does rubric-\allowbreak{}based automated judgment of task responses agree with domain-\allowbreak{}expert judgment of the same responses, measured by quadratic weighted kappa, relative to the agreement between experts themselves, and how does that agreement vary with item difficulty and across repeated runs? Expert--expert agreement provides the ceiling against which the automated judge is interpreted.

Both problems are evaluated in two stages, because the experts on whom they depend are engaged progressively (Section 4.2.1). In Stage A, during the present phase, the project team stands in for the experts. Team members play each free-text task and score the answers on the rubric, as experts will in Stage B; and the team constructs, for each task, a strong reference response and variants that each degrade one rubric criterion to a known level, together with probes for length, format, and ordering bias, following the perturbation-\allowbreak{}checklist method\cite{r73}. A panel of judging models from different families, none of which produced the synthetic responses, scores them independently and repeatedly\cite{r74}, and the advisor checks the team's scores and designed levels on a random sample of about a fifth. Stage A establishes whether the judge agrees with the team's scores, is sensitive to each criterion, robust to known biases, and consistent across runs; it cannot establish agreement with experts on real responses, and is not presented as doing so. In Stage B, on the completed product, volunteer participants complete tasks with informed consent, and their responses, together with generated scenarios, are rated independently and blind by three experts per pilot role. RP-1 is then answered by content validity indices, with I-CVI of at least .78 and S-CVI/Ave of at least .90\cite{r70}, by a paired comparison of seeded and unseeded generation using McNemar's test on acceptability, and by a fact-\allowbreak{}preservation target of 95\% of seed-\allowbreak{}required facts. RP-2 is answered against the three acceptance criteria of Williamson, Xi and Breyer\cite{r71}, with agreement among experts measured by Krippendorff's alpha\cite{r72}.

Three further questions arise in the project and are addressed, but they are not research problems under these conditions.

{\small\setstretch{1}\renewcommand{\arraystretch}{1.1}
\begin{longtable}{|>{\raggedright\arraybackslash}p{\dimexpr 0.349\textwidth-2\tabcolsep-\arrayrulewidth\relax}|>{\raggedright\arraybackslash}p{\dimexpr 0.486\textwidth-2\tabcolsep-\arrayrulewidth\relax}|>{\raggedright\arraybackslash}p{\dimexpr 0.159\textwidth-2\tabcolsep-\arrayrulewidth\relax}|}
\caption{Questions addressed in the project but not treated as research problems.} \\
\hline
\textbf{Question} & \textbf{Condition not met} & \textbf{Where it is addressed} \\ \hline
\endfirsthead
\hline
\textbf{Question} & \textbf{Condition not met} & \textbf{Where it is addressed} \\ \hline
\endhead
How should demonstrated competence determine access to adjacent occupations? & Resolved by a design decision based on occupational taxonomies; the project does not generate new knowledge about it & Design (Chapter 5) \\ \hline
Does the engagement design sustain exploration without displacing learning? & Answerable only through longitudinal user study, which the offline evaluation does not provide & Subsequent work \\ \hline
How should the expert reference set be constructed? & Concerns how evidence is gathered --- a matter of method rather than a problem & Evaluation (Chapter 7) \\ \hline
\end{longtable}
}

\subsubsection{Challenges and difficulties}

Generation fidelity is the principal technical difficulty. Occupational content is factual, and a scenario that is plausible but inaccurate is worse than no scenario, since it conveys a confident false impression of a real profession. Long-\allowbreak{}horizon consistency compounds this within multi-task levels.

Assessment validity is the principal methodological difficulty. Section 2.5 indicates that agreement figures for automated judgment are frequently overstated and that performance degrades unevenly with difficulty; obtaining defensible figures requires a reference set that is stratified and expert-\allowbreak{}authored, which is expensive to produce.

Expert availability constrains breadth. Each role requires three to four senior practitioners for authoring and review, whose time must be arranged through partner organisations; in Stage 1, an estimated seven hours per expert covers scenario ratings, response ratings, and option keying for a pilot role. The three pilot roles are back-end engineering, DevOps and cloud engineering, and business analysis, the last chosen so that assessment is also examined on work whose output is requirements and communication with stakeholders rather than code. Because experts are engaged progressively, the present phase relies on provisional seeds prepared from public sources, and occupational accuracy can be claimed only once the Stage 1 panel has reviewed them. This is the reason for restricting the first phase to a single field, and for treating expert authoring as a first-class activity of the system rather than an administrative convenience.

Cold start affects the aggregate occupational profiles, which require a population of completions before they are informative. Until that population exists, profiles are absent or unreliable, and the interface must represent this honestly rather than display an unfounded figure.

Evaluation without deployment limits what the project can establish. Offline benchmarking substantiates generation quality and assessment accuracy but cannot establish engagement retention or effects on career decisions, which require longitudinal study. The claims are bounded accordingly, consistent with Section 2.8.

Operating without user revenue makes cost a design constraint rather than a commercial concern. Inference cost per session must remain within what grant funding can sustain, which constrains model choice, response length, and caching strategy.

Interaction modality limits what the platform can represent, and the limit differs by competency. Hamstra et al.\cite{r98} argue that the realism of a simulation should be understood not as physical resemblance but as functional task alignment --- the correspondence between what a task demands and what the real situation demands --- noting that physical fidelity has been found largely independent of educational effectiveness. On this view the relevant question is not whether a screen-\allowbreak{}based task looks like a workplace but whether it demands what the work demands. Table 4.5 compares the competencies a JobQuest task can involve with their counterparts in a live work situation, rating alignment on four levels: high, moderate, partial, and low.

{\small\setstretch{1}\renewcommand{\arraystretch}{1.1}
\begin{longtable}{|>{\raggedright\arraybackslash}p{\dimexpr 0.264\textwidth-2\tabcolsep-\arrayrulewidth\relax}|>{\raggedright\arraybackslash}p{\dimexpr 0.350\textwidth-2\tabcolsep-\arrayrulewidth\relax}|>{\raggedright\arraybackslash}p{\dimexpr 0.222\textwidth-2\tabcolsep-\arrayrulewidth\relax}|>{\raggedright\arraybackslash}p{\dimexpr 0.159\textwidth-2\tabcolsep-\arrayrulewidth\relax}|}
\caption{Alignment of competencies between screen-\allowbreak{}based tasks and live work situations.} \\
\hline
\textbf{Competency} & \textbf{In a JobQuest task} & \textbf{In a live work situation} & \textbf{Alignment} \\ \hline
\endfirsthead
\hline
\textbf{Competency} & \textbf{In a JobQuest task} & \textbf{In a live work situation} & \textbf{Alignment} \\ \hline
\endhead
Technical analysis and problem-\allowbreak{}solving, such as diagnosing a failing service & Performed through the same text and code interfaces used in practice & Largely screen-\allowbreak{}based in IT roles & High \\ \hline
Written communication, such as incident reports, documentation, and messages to colleagues & Produced and assessed as text & Produced as text & High \\ \hline
Prioritisation and decision-\allowbreak{}making under constraints & Decisions made with consequences represented in the scenario & The same decisions, with real consequences for users and colleagues & Moderate: decision structure aligned, stakes attenuated \\ \hline
Adaptability to changing requirements & Changes introduced as scenario events & Changes arrive unpredictably and compound & Moderate \\ \hline
Collaboration and handling of stakeholders & Text interaction with AI-driven colleagues and clients & Real people with their own goals, relationships, and history & Partial: structure aligned, social dynamics simplified \\ \hline
Working under time pressure and regulating stress & Time limits can be simulated, with low personal stakes & Sustained pressure with genuine stakes & Low \\ \hline
Oral communication and presentation & Represented only through what is written; delivery, tone, and listening are not exercised & Central in many roles & Low \\ \hline
Interviews & The content of answers can be rehearsed in text; live delivery under observation is not exercised & The gateway to most roles & Low \\ \hline
Physical and manual skill & Procedures can be described and sequenced; manual execution is not exercised & Central in some occupations & Low \\ \hline
\end{longtable}
}

Two consequences follow. Alignment is highest for the competencies at the core of information technology work, which is itself largely screen-\allowbreak{}mediated --- a further reason for the first phase's coverage. Conversely, competencies that depend on live interpersonal dynamics, genuine stakes, speech, or physical performance are attenuated: oral communication, interviewing, and physical skill are represented only through their written or procedural content, and alignment for them is low. Both the assessment and the self-\allowbreak{}efficacy the platform builds (Section 2.4) are therefore evidence about the well-\allowbreak{}aligned competencies only, and results must not be presented as indicating readiness for the others. Voice interaction could extend coverage to oral communication; genuine stakes and physical skill lie outside what any screen-\allowbreak{}based simulation can provide.

Localisation presents a further difficulty: occupational taxonomies are authored in English for other labour markets, and their occupational structures do not map cleanly onto Vietnamese practice, requiring expert adaptation rather than translation.

\subsection{User's Requirements}

Requirements are expressed first as user requirements per group, each with a priority (Must, Should, Could). These are then analysed into functional, non-\allowbreak{}functional, and data requirements. Each of these carries two verification methods: a business verification, by which a stakeholder --- an explorer, a domain expert, a content administrator, or a system administrator --- confirms that the requirement meets its purpose, through acceptance tests, walkthroughs, or expert review; and a technical verification, by which the development team confirms that the implementation behaves correctly, through automated, integration, boundary, load, or fault-\allowbreak{}injection tests. The requirements are aligned with the use case diagram UC-00, the activity diagrams AD-01a to AD-05e, and the state machine SM-01 of Chapter 5. The business verification and the technical verification of each functional, non-\allowbreak{}functional, and data requirement are given in Appendix 3.

\subsubsection{User requirements}

The user requirements are 20 for explorers (UR-E01 to UR-E20), 9 for domain experts (UR-D01 to UR-D09), 8 for content administrators (UR-C01 to UR-C08), and 4 for system administrators (UR-S01 to UR-S04). Each is written as a user story with a priority of Must, Should, or Could, and all are listed in Appendix 3. The functional, non-\allowbreak{}functional, and data requirements below are derived from them.

\subsubsection{Functional requirements}

{\small\setstretch{1}\renewcommand{\arraystretch}{1.1}
\begin{longtable}{|>{\raggedright\arraybackslash}p{\dimexpr 0.103\textwidth-2\tabcolsep-\arrayrulewidth\relax}|>{\raggedright\arraybackslash}p{\dimexpr 0.699\textwidth-2\tabcolsep-\arrayrulewidth\relax}|>{\raggedright\arraybackslash}p{\dimexpr 0.193\textwidth-2\tabcolsep-\arrayrulewidth\relax}|}
\caption{Functional requirements.} \\
\hline
\textbf{ID} & \textbf{Functional requirement} & \textbf{Group} \\ \hline
\endfirsthead
\hline
\textbf{ID} & \textbf{Functional requirement} & \textbf{Group} \\ \hline
\endhead
FR-01 & Register, sign in with a password or Google, and recover the account & Explorer \\ \hline
FR-02 & Edit, sign out of, and delete the account & Explorer \\ \hline
FR-03 & Run an optional, retakable RIASEC orientation quiz suggesting occupations\cite{r41} & Explorer \\ \hline
FR-04 & Show the explorer's portrait for their RIASEC profile, including the social index & Explorer \\ \hline
FR-05 & Browse occupations on the Galaxy Map and show each profile with its downsides and level descriptions & Explorer \\ \hline
FR-06 & Open the entry level of every published occupation to every explorer & Explorer \\ \hline
FR-07 & Let explorers start an adjacent occupation at a higher level when overlap conditions are met & Explorer \\ \hline
FR-08 & Enrol in levels, choose scenarios, award credit, and unlock the next level & Explorer \\ \hline
FR-09 & Offer the role, team, and context before each scenario & Explorer \\ \hline
FR-10 & Support choice, ordering, prioritising, and free-text answers & Explorer \\ \hline
FR-11 & Score each free-text answer, including every follow-up answer, with anchor points backed by quotations & Explorer \\ \hline
FR-12 & Answer follow-up questions generated by AI on free-text answers & Explorer \\ \hline
FR-13 & Offer hints when the explorer is stuck & Explorer \\ \hline
FR-14 & Return result, feedback, and guidance separately & Explorer \\ \hline
FR-15 & Record evidenced expertise and soft-skill competencies, capped per skill per scenario & Explorer \\ \hline
FR-16 & Offer optional side quests for bonus points & Explorer \\ \hline
FR-17 & Show the career path of each occupation level & Explorer \\ \hline
FR-18 & Show a private record of attempts, expertise and soft-skill points, awards, and past scenario feedback & Explorer \\ \hline
FR-19 & Collect occupation ratings and feedback & Explorer \\ \hline
FR-20 & Send app feedback and review it & Explorer, System administrator \\ \hline
FR-21 & Manage the occupation profile: levels, level descriptions, future roles, downsides, salary, generation rules, competencies, adjacency & Domain expert \\ \hline
FR-22 & Author scenario context, allowed task types, option scores, rubric, hints, random events, and side quests & Domain expert \\ \hline
FR-23 & Generate scenario content with AI from the author's context, keeping its generation history & Domain expert \\ \hline
FR-24 & Add reference answers with expert scores & Domain expert \\ \hline
FR-25 & Review another expert's scenario against the review criteria & Domain expert \\ \hline
FR-26 & Play-test the draft game before publication & Domain expert \\ \hline
FR-27 & Revalidate a published scenario on request & Domain expert \\ \hline
FR-28 & Report disagreements between AI and expert grading & Domain expert \\ \hline
FR-29 & Rate responses and scenarios blind & Domain expert \\ \hline
FR-30 & Create and manage own occupations & Content administrator \\ \hline
FR-31 & Import O*NET and ESCO reference data & Content administrator \\ \hline
FR-32 & Define the scenario review criteria & Content administrator \\ \hline
FR-33 & Open one active scenario per task & Content administrator \\ \hline
FR-34 & Assign a domain expert as reviewer, with a deadline, and reassign overdue reviews & Content administrator \\ \hline
FR-35 & Choose the scenario presentation with AI: task types from the domain expert's allowed set, characters, and scenes & Content administrator \\ \hline
FR-36 & Maintain a bias-\allowbreak{}checked asset library; AI generates missing assets & Content administrator \\ \hline
FR-37 & Publish a scenario after review and play-test & Content administrator \\ \hline
FR-38 & Delete a published scenario by retiring it & Content administrator \\ \hline
FR-39 & Monitor own occupations: statistics and coverage & Content administrator \\ \hline
FR-40 & Assign and reassign one content administrator and domain experts to each occupation & System administrator \\ \hline
FR-41 & Manage all user accounts, roles, and role-based access & System administrator \\ \hline
FR-42 & Configure the generation model and the three judge models separately & System administrator \\ \hline
FR-43 & Monitor usage, cost, and errors, and re-run failed AI jobs & System administrator \\ \hline
\end{longtable}
}

\subsubsection{Non-\allowbreak{}functional requirements}

{\small\setstretch{1}\renewcommand{\arraystretch}{1.1}
\begin{longtable}{|>{\raggedright\arraybackslash}p{\dimexpr 0.114\textwidth-2\tabcolsep-\arrayrulewidth\relax}|>{\raggedright\arraybackslash}p{\dimexpr 0.665\textwidth-2\tabcolsep-\arrayrulewidth\relax}|>{\raggedright\arraybackslash}p{\dimexpr 0.216\textwidth-2\tabcolsep-\arrayrulewidth\relax}|}
\caption{Non-\allowbreak{}functional requirements.} \\
\hline
\textbf{ID} & \textbf{Non-\allowbreak{}functional requirement} & \textbf{Category} \\ \hline
\endfirsthead
\hline
\textbf{ID} & \textbf{Non-\allowbreak{}functional requirement} & \textbf{Category} \\ \hline
\endhead
NFR-01 & Follow-up generation within 5 s (95th percentile); progress shown after 2 s & Performance \\ \hline
NFR-02 & Free-text assessment within 8 s (95th percentile) & Performance \\ \hline
NFR-03 & Progress survives interruption & Reliability \\ \hline
NFR-04 & Failover to another model when a model fails; no progress lost and user notified within 10 s when no model is available & Reliability \\ \hline
NFR-05 & Minimal, encrypted, deletable personal data; only de-\allowbreak{}identified answers sent to AI providers & Security and privacy \\ \hline
NFR-06 & Suppress aggregates below a minimum population & Security and privacy \\ \hline
NFR-07 & Disclose AI generation and AI assessment & Transparency \\ \hline
NFR-08 & Vietnamese-\allowbreak{}language interface and content & Usability \\ \hline
NFR-09 & At least 80\% of readers state the task objective correctly & Usability \\ \hline
NFR-10 & New occupations need no code change & Maintainability \\ \hline
NFR-11 & Session cost stays within the configured ceiling & Cost \\ \hline
NFR-12 & Model provider replaceable without content changes & Portability \\ \hline
NFR-13 & Works in a mobile browser on a slow connection & Compatibility \\ \hline
NFR-14 & AI grading validated against expert judgment & Fairness \\ \hline
NFR-15 & Published scenarios match approved content; AI content during play stays within expert-set limits & Accuracy \\ \hline
NFR-16 & Scenario status changes logged; published content immutable & Auditability \\ \hline
NFR-17 & Passwords hashed, all traffic encrypted, idle sessions expire, AI prompts protected against injection & Security and privacy \\ \hline
NFR-18 & Explorer service available within the agreed operating target & Availability \\ \hline
NFR-19 & Explorer screens meet WCAG 2.1 level AA & Accessibility \\ \hline
\end{longtable}
}

Design constraints are recorded separately from business rules, since they arise from implementation choices rather than from the domain. Each derives from at least one non-\allowbreak{}functional requirement: the requirement states what must be achieved, and the constraint fixes or prohibits a design choice so that it can be.

{\small\setstretch{1}\renewcommand{\arraystretch}{1.1}
\begin{longtable}{|>{\raggedright\arraybackslash}p{\dimexpr 0.103\textwidth-2\tabcolsep-\arrayrulewidth\relax}|>{\raggedright\arraybackslash}p{\dimexpr 0.477\textwidth-2\tabcolsep-\arrayrulewidth\relax}|>{\raggedright\arraybackslash}p{\dimexpr 0.126\textwidth-2\tabcolsep-\arrayrulewidth\relax}|>{\raggedright\arraybackslash}p{\dimexpr 0.289\textwidth-2\tabcolsep-\arrayrulewidth\relax}|}
\caption{Design constraints.} \\
\hline
\textbf{ID} & \textbf{Design constraint} & \textbf{Derived from} & \textbf{Rationale} \\ \hline
\endfirsthead
\hline
\textbf{ID} & \textbf{Design constraint} & \textbf{Derived from} & \textbf{Rationale} \\ \hline
\endhead
DC-01 & The three judge models and the generation model shall come from four different model families. & NFR-14 & Reduces correlated failure modes and self-\allowbreak{}preference bias, thereby improving the independence of automated evaluation (Section 2.5) \\ \hline
DC-02 & The system shall support replacing the underlying model provider without requiring changes to the core application logic. & NFR-12, NFR-04 & Dependence on third-party model providers; portability as funding conditions change \\ \hline
DC-03 & Personal data is processed and stored in compliance with Vietnam's Law on Personal Data Protection and its implementing decree, including the provisions for cross-\allowbreak{}border transfer. & NFR-05, NFR-06 & Legal obligation (Sections 2.7, 4.1.4) \\ \hline
DC-05 & The platform is delivered as a web application usable through standard browsers without installation. & NFR-13 & Universal access (Section 4.1.5) \\ \hline
\end{longtable}
}

\subsubsection{Data requirements}

{\small\setstretch{1}\renewcommand{\arraystretch}{1.1}
\begin{longtable}{|>{\raggedright\arraybackslash}p{\dimexpr 0.103\textwidth-2\tabcolsep-\arrayrulewidth\relax}|>{\raggedright\arraybackslash}p{\dimexpr 0.587\textwidth-2\tabcolsep-\arrayrulewidth\relax}|>{\raggedright\arraybackslash}p{\dimexpr 0.305\textwidth-2\tabcolsep-\arrayrulewidth\relax}|}
\caption{Data requirements.} \\
\hline
\textbf{ID} & \textbf{Data requirement} & \textbf{Source} \\ \hline
\endfirsthead
\hline
\textbf{ID} & \textbf{Data requirement} & \textbf{Source} \\ \hline
\endhead
DR-01 & Occupational reference data (occupations, tasks, skills) & O*NET and ESCO, imported by the content administrator and adapted for Vietnamese practice by domain experts \\ \hline
DR-02 & Competency taxonomy with observable indicators & Derived from DR-01 under expert adaptation \\ \hline
DR-03 & Occupational adjacency with overlap measures & Computed from DR-01/DR-02 \\ \hline
DR-04 & Scenario content: context, activities, option scores, rubric, hints, random events, side quests & Author domain expert with AI assistance \\ \hline
DR-05 & Generated scenarios with provenance & System \\ \hline
DR-06 & User progress: enrolled levels, competence credit capped per scenario, side-quest points, play history, end-\allowbreak{}of-\allowbreak{}scenario feedback, awards & System \\ \hline
DR-07 & Assessment records, including each judge's score and the combined grade per answer, hidden feedback notes, hint use, and explorer flags & System and AI judges, with flags from explorers \\ \hline
DR-08 & Five-\allowbreak{}attribute occupation ratings (project-\allowbreak{}defined) & Explorers \\ \hline
DR-09 & Expert reference set stratified by difficulty & Domain experts \\ \hline
DR-10 & Operational telemetry & System \\ \hline
DR-11 & Play-test session flag & System \\ \hline
DR-12 & Quiz answers, RIASEC profile with social index, and the fixed portrait texts (one per RIASEC type) & Explorers; portrait texts written in advance by the project team \\ \hline
DR-13 & Scenario lifecycle record, including author and reviewer assignments with review deadlines, review decisions, comments, and revalidation results & System, from staff actions \\ \hline
DR-14 & Presentation data and asset library & Content administrator with AI \\ \hline
DR-15 & Occupation catalogue (domains and occupations) and occupation profiles (level descriptions, future roles, downsides) & Content administrator (catalogue) and domain experts (profiles) \\ \hline
DR-16 & User accounts, linked Google identities, roles, assignments, and activity log & System administrator; activity log from the system \\ \hline
DR-17 & Platform settings & System administrator \\ \hline
DR-18 & App feedback records & Explorers \\ \hline
DR-19 & Scenario review criteria & Content administrator \\ \hline
\end{longtable}
}

Personal data within DR-06, DR-07, DR-08, DR-11, DR-12, DR-16, and DR-18 is subject to the policies in Section 4.1.4. DR-09 contains no personal data and is retained as a research asset.

\subsection{Summary}

This chapter specified the proposed system against the domain established in Chapters 1 to 3. The business context identified three operating processes --- exploration, content lifecycle, and evaluation --- and twelve business rules, several of which derive from the evidence reviewed in Chapter 2 rather than from product preference: accessible entry levels, competence-\allowbreak{}gated progression, taxonomy-\allowbreak{}derived adjacency, mandatory expert validation, and formative-\allowbreak{}only assessment. It set out a public-\allowbreak{}benefit business model in which all content is free to users, operation is funded through grants, acknowledged sponsorship, donations, and an intended partnership with the Ministry of Education and Training, and sponsors are excluded from content so that the candour of occupational representation is protected, together with a five-stage roadmap in which each stage is gated by evidence, and a staged organisation of five operating units under a steering committee, with content governed by occupational domain. The system description identified four user groups --- explorers, domain experts, content administrators, and system administrators --- two research problems in AI applied to education --- grounded generation and the validity of stealth assessment, each stated with the comparison that makes its answer knowledge --- and the principal difficulties, the most consequential being generation fidelity, assessment validity, expert availability, cost sustainability without revenue, and the uneven alignment between screen-\allowbreak{}based tasks and live work. The requirements analysis produced forty-three prioritised user requirements, from which forty-four functional, nineteen non-\allowbreak{}functional, and nineteen data requirements were derived, each paired with a business and a technical verification method, together with four design constraints, each derived from a non-\allowbreak{}functional requirement and recorded separately from the business rules. Chapter 5 models the processes and interactions these requirements imply.
